\section{Joint contractions of the free source inputs}
\label{sec:peps_source_block_contraction}

Fix ordered source slots with halfspaces $U_e,V_e$ and joint register
space $\mathcal S$, as in Definition~\ref{def:peps_source_slots}.
All source and physical input coordinates in this section remain free.
This gives the contractions used in
\cite[proof of Theorem~5.2]{OpenAI2026PolynomialPEPS},
\texttt{04-compression.tex}, lines 383--468.

% The operators below come from the specified allowed compositions.
% Constructing the affected and exterior compositions from circuit chronology,
% and deriving the corrected-subset error estimate, remain separate obligations.
% Individual prepared-matrix contraction bounds are not used to infer a
% contraction bound on their combined matrix.

\begin{definition}[Bases of source and input registers]
    \label{def:peps_source_slot_basis}
    \lean{TNLean.PEPS.PairEffect.SourceInventory.slotBasis}
    \lean{TNLean.PEPS.PairEffect.SourceInventory.slotBasis_apply}
    \lean{TNLean.PEPS.PairEffect.SourceInventory.preparedInputBasis}
    \lean{TNLean.PEPS.PairEffect.SourceInventory.preparedInputBasis_apply}
    \leanok
    \uses{def:peps_joint_source_vector}
    Choose finite orthonormal bases $(u_{e,r})_{r\in I_e}$ of $U_e$
    and $(v_{e,s})_{s\in J_e}$ of $V_e$. Let
    $\mathcal A=\prod_e(I_e\times J_e)$.
    The orthonormal basis of $\mathcal S$ indexed by $a\in\mathcal A$
    consists of
    \begin{align}
        s_a & =\Phi\bigl((u_{e,a_e^{(1)}}\otimes
        v_{e,a_e^{(2)}})_e\bigr).
        \label{eq:peps_source_slot_basis}
    \end{align}
    The empty source list has its single basis vector $1\in\C$.
    If $(f_j)_{j\in N}$ is an orthonormal basis of the original
    input memory $\mathcal H$, the resulting orthonormal basis of
    $\mathcal S\otimes\mathcal H$ is indexed by $(a,j)$ and satisfies
    \begin{align}
        s_a\otimes f_j
         & =Q_{(u_{e,a_e^{(1)}}\otimes v_{e,a_e^{(2)}})_e,\mathcal H}f_j.
        \label{eq:peps_source_input_basis}
    \end{align}
    Both formulas use the canonical associativity identifications.
\end{definition}

\begin{proof}\leanok
    \uses{thm:peps_source_preparation_sums}
    Induct on the ordered slots. Tensor the bases of the first two
    halfspaces with the basis already constructed for the remaining
    slots, and identify a full assignment with its first component
    and its restriction to the remaining slots. Tensor products and
    isometric identifications preserve orthonormal bases. This gives
    (\ref{eq:peps_source_slot_basis}); tensoring with the input basis
    and applying the preparation identity gives
    (\ref{eq:peps_source_input_basis}).
\end{proof}

\begin{definition}[The complete free-input matrix]
    \label{def:peps_free_source_matrix}
    \lean{TNLean.PEPS.PairEffect.Word.freeSourceMatrix}
    \lean{TNLean.PEPS.PairEffect.Word.freeSourceMatrix_apply}
    \leanok
    \uses{def:peps_source_slot_basis,def:peps_prepared_matrix,def:peps_prepared_density_coefficient}
    Let $T:\mathcal S\otimes\mathcal H\to\mathcal K$ be the
    operator of a composition, and choose an orthonormal output
    basis $(g_k)_{k\in M}$. Its complete matrix $A_T$ has row index
    $k\in M$ and column index $(a,j)\in\mathcal A\times N$:
    \begin{align}
        A_T(k,(a,j)) & =\langle g_k,T(s_a\otimes f_j)\rangle
        =K_T(a)_{kj}.
        \label{eq:peps_free_source_matrix}
    \end{align}
    Here $K_T(a)$ is the actual prepared matrix of
    Definition~\ref{def:peps_prepared_density_coefficient}.
\end{definition}

\begin{theorem}[Joint contraction of all source inputs]
    \label{thm:peps_free_source_contraction}
    \lean{TNLean.PEPS.PairEffect.Word.norm_freeSourceMatrix_le_one}
    \leanok
    \uses{def:peps_free_source_matrix}
    If $T$ is an allowed composition, then $\|A_T\|\le1$ in the
    Euclidean operator norm. This bounds its action on arbitrary
    superpositions of the source and physical input coordinates,
    independently of the number of slots and their dimensions.
\end{theorem}

\begin{proof}\leanok
    \uses{def:peps_party_layout,def:peps_source_slot_basis}
    The complete matrix is the matrix of the same operator $T$ in
    orthonormal input and output bases. Coordinate maps preserve
    norms, so $\|A_T\|=\|T\|\le1$.
\end{proof}

\begin{theorem}[The bra-adjoint--ket contraction]
    \label{thm:peps_free_source_overlap}
    \lean{TNLean.PEPS.PairEffect.Word.freeSourceMatrix_conjTranspose_mul_apply}
    \lean{TNLean.PEPS.PairEffect.Word.norm_freeSourceMatrix_conjTranspose_mul_le_one}
    \leanok
    \uses{def:peps_free_source_matrix}
    Let $T$ and $T'$ be compositions with a common output
    memory and output basis. Their source lists, source spaces,
    source bases, and physical input memories and bases may differ.
    Write $A=A_T$ and $B=A_{T'}$, and let $a,b$ denote the respective
    source assignments. Then the rectangular matrix $O=B^\dagger A$
    satisfies
    \begin{align}
        O_{(b,j),(a,i)} & =(K_{T'}(b)^\dagger K_T(a))_{ji}.
        \label{eq:peps_free_source_overlap}
    \end{align}
    Thus the bra source assignment and its physical input index form
    the row index of $O$. If both compositions are allowed, then
    $\|O\|\le1$.
\end{theorem}

\begin{proof}\leanok
    \uses{thm:peps_free_source_contraction}
    Multiplying the matrices in (\ref{eq:peps_free_source_matrix})
    gives (\ref{eq:peps_free_source_overlap}). The norm bound follows
    from $\|B^\dagger A\|\le\|B^\dagger\|\|A\|=\|B\|\|A\|\le1$.
\end{proof}

\begin{theorem}[Tracing a prepared density coefficient]
    \label{thm:peps_free_source_trace}
    \lean{TNLean.PEPS.PairEffect.Word.trace_preparedMatrix_mul_conjTranspose_eq}
    \lean{TNLean.PEPS.PairEffect.Word.trace_preparedDensityCoefficient_eq_freeSourceMatrix}
    \leanok
    \uses{def:peps_free_source_matrix,def:peps_prepared_density_coefficient}
    With the ket and bra input coordinates as above, let
    $\rho\in\C^{N\times N'}$ be any rectangular matrix. No
    contraction assumption on $T$ or $T'$ is needed for the identity
    \begin{align}
        \tr\bigl(K_T(a)\rho K_{T'}(b)^\dagger\bigr)
         & =\sum_{i\in N}\sum_{j\in N'}
        \rho_{ij}\,O_{(b,j),(a,i)}.
        \label{eq:peps_free_source_trace}
    \end{align}
    When the two branches have common source spaces and physical
    input coordinates, this is the trace of
    $C_{T,T'}(\rho;a,b)$ from
    Definition~\ref{def:peps_prepared_density_coefficient}, with
    independent endpoint bases on the two sides.
\end{theorem}

\begin{proof}\leanok
    \uses{thm:peps_free_source_overlap}
    Cyclicity of the rectangular trace gives
    $\tr(K_T(a)\rho K_{T'}(b)^\dagger)
        =\tr(\rho K_{T'}(b)^\dagger K_T(a))$.
    Expand the trace and substitute the block identity
    (\ref{eq:peps_free_source_overlap}).
\end{proof}

\begin{definition}[Pair labels of the actual slots]
    \label{def:peps_source_pair_index}
    \lean{TNLean.PEPS.PairEffect.SourceInventory.pairIndexEquiv}
    \lean{TNLean.PEPS.PairEffect.SourceInventory.pairIndexEquiv_apply}
    \leanok
    \uses{def:peps_source_inventory}
    Let $R$ be a source list whose unordered pair labels are distinct
    and contain every unordered pair of distinct parties exactly once.
    Sending a slot to the unordered pair of its actual source is a
    bijection from its slot indices to
    $\{\{p,q\}:p,q\in P,\ p\ne q\}$.
    This identification depends only on the parties of each source.
\end{definition}

\begin{theorem}[Number of pair-source slots]
    \label{thm:peps_source_pair_count}
    \lean{TNLean.PEPS.PairEffect.SourceInventory.length_eq_choose}
    \lean{TNLean.PEPS.PairEffect.SourceInventory.length_le_card_sq}
    \leanok
    \uses{def:peps_source_pair_index}
    If $P$ is finite and $R$ has the pair coverage and distinctness
    in Definition~\ref{def:peps_source_pair_index}, then
    $|R|=\binom{|P|}{2}\le|P|^2$.
    This includes empty and singleton party sets.
    These are the pair slots used in
    \cite[proof of Theorem~5.2]{OpenAI2026PolynomialPEPS},
    \texttt{04-compression.tex}, lines 565--585.
\end{theorem}

\begin{proof}\leanok
    \uses{def:peps_source_pair_index}
    Count the unordered pairs of distinct elements under the stated
    bijection. Their number is $\binom{|P|}{2}$, which is at most
    $|P|^2$.
\end{proof}
\begin{definition}[Local maps on source slots]
    \label{def:peps_source_slot_maps}
    \lean{TNLean.PEPS.PairEffect.Word.mapSourceSlots}
    \leanok
    \uses{def:peps_source_slots,def:peps_party_layout}
    For each slot $e$, let $L_e:U_e\to U'_e$ and
    $M_e:V_e\to V'_e$ be continuous linear maps. Apply these maps
    to the corresponding source registers and leave the original
    input memory $\mathcal H$ unchanged. Denote the resulting
    composition by $F_{L,M}$. This composition depends on the maps
    and the ordered slots, but not on the prepared source vectors.
\end{definition}

\begin{theorem}[Action of the local source maps]
    \label{thm:peps_source_slot_maps}
    \lean{TNLean.PEPS.PairEffect.Word.mapSourceSlots_spec}
    \lean{TNLean.PEPS.PairEffect.Word.eval_mapSourceSlots_prepareSlots}
    \leanok
    \uses{def:peps_source_slot_maps,def:peps_source_slots}
    For every family $\eta_e\in U_e\otimes V_e$,
    \begin{align}
        F_{L,M}Q_{\eta,\mathcal H}
         & =Q_{((L_e\otimes M_e)\eta_e)_e,\mathcal H}.
        \label{eq:peps_source_slot_maps}
    \end{align}
    If $\|L_e\|,\|M_e\|\le1$ for every slot, the composition
    $F_{L,M}$ is allowed and contains no source preparations.
    The identity requires neither these norm bounds nor normalization
    of the vectors $\eta_e$.
\end{theorem}

\begin{proof}\leanok
    \uses{thm:peps_map_source_pair,thm:peps_source_preparation}
    Induct on the ordered source slots. Apply the two endpoint maps
    of the first source, and then apply the construction to the
    remaining sources while leaving those first two registers fixed.
    Moving the remaining preparation past these fixed registers gives
    (\ref{eq:peps_source_slot_maps}). Each constituent is a local map,
    and its allowedness follows from the corresponding endpoint norm
    bound.
\end{proof}

\begin{definition}[The complete matrix in source frames]
    \label{def:peps_source_frame_matrix}
    \lean{TNLean.PEPS.PairEffect.Word.sourceFrameMatrix}
    \lean{TNLean.PEPS.PairEffect.Word.sourceFrameMatrix_eq_freeSourceMatrix}
    \leanok
    \uses{def:peps_source_slot_maps,def:peps_free_source_matrix}
    Let $T:\mathcal S'\otimes\mathcal H\to\mathcal K$ be a
    composition on the target halfspaces $U'_e,V'_e$. Use the finite
    orthonormal bases of $U_e,V_e$ and the physical input and output
    bases already specified. Define
    \begin{align}
        A_T^{L,M}(k,(a,j))
         & =\left\langle g_k,
        TQ_{(L_eu_{e,a_e^{(1)}}\otimes
                M_ev_{e,a_e^{(2)}})_e,\mathcal H}f_j\right\rangle.
        \label{eq:peps_source_frame_matrix}
    \end{align}
    This is the complete free-input matrix of $T F_{L,M}$:
    \begin{align}
        A_T^{L,M} & =A_{TF_{L,M}}.
        \label{eq:peps_source_frame_matrix_composition}
    \end{align}
    The images of the endpoint maps need not span the target
    halfspaces.
\end{definition}

\begin{proof}\leanok
    \uses{thm:peps_source_slot_maps}
    Evaluate (\ref{eq:peps_source_slot_maps}) on each elementary
    endpoint-basis preparation and each physical input basis vector.
    The resulting columns agree with
    (\ref{eq:peps_source_frame_matrix}).
\end{proof}

\begin{theorem}[Joint contraction in source frames]
    \label{thm:peps_source_frame_contraction}
    \lean{TNLean.PEPS.PairEffect.Word.norm_sourceFrameMatrix_le_one}
    \leanok
    \uses{def:peps_source_frame_matrix}
    Assume that $T$ is allowed and that each $L_e$ and $M_e$ is a
    contraction. Then $\|A_T^{L,M}\|\le1$ in the Euclidean operator norm.
    In particular, this applies to isometric inclusions of proper
    Schmidt subspaces. No orthonormal basis of the target halfspaces
    is required.
\end{theorem}

\begin{proof}\leanok
    \uses{thm:peps_source_slot_maps,thm:peps_free_source_contraction}
    Each $L_e$ and $M_e$ is a contraction, so the composition
    $TF_{L,M}$ is allowed. By (\ref{eq:peps_source_frame_matrix_composition})
    and Theorem~\ref{thm:peps_free_source_contraction},
    \begin{align*}
        \|A_T^{L,M}\| & =\|A_{TF_{L,M}}\|\le1.
    \end{align*}
\end{proof}
